\begin{ouverture}[courbe]{10}{Statistiques et représentations}{Partie 3 · Données et hasard}
{Effectifs · Fréquences · Moyenne · Médiane · Étendue}
\objectif{Organiser et représenter une série.}
\objectif{Calculer et interpréter des indicateurs.}
\objectif{Comparer sans confondre moyenne et dispersion.}
\prerequis{Fractions et pourcentages.}
\end{ouverture}
\pagestyle{otMain}\essentiel\fontsize{11}{15}\selectfont
\essSection{1}{Décrire une série}
\begin{definition}L’effectif total est le nombre d’observations. La fréquence d’une valeur est son effectif divisé par l’effectif total. La somme des fréquences vaut 1, soit 100 \%.\end{definition}
\begin{propriete}La moyenne est la somme des valeurs divisée par leur nombre. Avec un tableau, multiplier chaque valeur par son effectif, additionner les produits puis diviser par l’effectif total.\end{propriete}
\begin{definition}Une médiane partage une série ordonnée en deux groupes : au moins la moitié des valeurs sont inférieures ou égales à elle et au moins la moitié sont supérieures ou égales à elle. Pour un effectif pair, ce livre retient la moyenne des deux valeurs centrales.\end{definition}
\begin{definition}L’étendue est la différence entre la plus grande et la plus petite valeur. Elle mesure une dispersion, sans décrire la répartition entre les extrêmes.\end{definition}
\begin{exemple}Pour $2;3;3;8$, la moyenne vaut 4, la médiane vaut 3 et l’étendue vaut 6. Ces trois nombres répondent à des questions différentes.\end{exemple}
\clearpage
\essSection{2}{Choisir et lire un graphique}
\begin{savoirfaire}[Représenter]Un diagramme en bâtons convient à des valeurs distinctes ; un diagramme circulaire partage un total, avec un angle égal à $360^\circ$ multiplié par la fréquence. Un histogramme représente des données regroupées en classes.\end{savoirfaire}
\begin{propriete}Dans les histogrammes de ce chapitre, les classes ont la même largeur. Les hauteurs peuvent donc représenter directement leurs effectifs. Les rectangles sont accolés. Une classe $[10;20[$ contient 10 et exclut 20.\end{propriete}
\begin{attention}Un regroupement en classes fait perdre les valeurs exactes. On peut repérer une classe contenant la médiane, mais pas forcément la médiane exacte. Les centres de classes donnent seulement une estimation de la moyenne.\end{attention}
\begin{savoirfaire}[Comparer]Préciser les unités, la taille des groupes, les axes et leur origine. Une moyenne identique n’impose ni une médiane identique ni une même dispersion. Une observation extrême peut déplacer fortement la moyenne.\end{savoirfaire}
\clearpage\rubrique{Découverte}
\begin{blocLivre}{ex-10-1}\begin{exercice}[1]{Compter}
Pour la série $2;4;4;5;7;8$, donner l’effectif total et l’effectif de 4.
\end{exercice}\end{blocLivre}
\begin{blocLivre}{ex-10-2}\begin{exercice}[1]{Une fréquence}
Dans un groupe de 30 élèves, 12 viennent à vélo. Donner la fréquence et le pourcentage correspondants.
\end{exercice}\end{blocLivre}
\begin{blocLivre}{ex-10-3}\begin{exercice}[1]{Une moyenne simple}
Calculer la moyenne de $7;9;10;14$.
\end{exercice}\end{blocLivre}
\begin{blocLivre}{ex-10-4}\begin{exercice}[1]{Ordonner avant la médiane}
Déterminer la médiane de $13;4;9;7;11$.
\end{exercice}\end{blocLivre}
\begin{blocLivre}{ex-10-5}\begin{exercice}[1]{Mesurer l’étendue}
Les températures relevées sont $-3;2;5;-1;8$ degrés Celsius. Donner l’étendue.
\end{exercice}\end{blocLivre}

\clearpage\rubrique{Entraînement}
\begin{blocLivre}{ex-10-6}\begin{exercice}[2]{Un tableau complet}
Calculer l’effectif, la moyenne, la médiane et l’étendue des notes.\begin{center}\renewcommand{\arraystretch}{1.35}\begin{tabular}{c|ccccc}\toprule
Note & $6$ & $8$ & $10$ & $12$ & $16$\\\midrule
Effectif & $2$ & $3$ & $6$ & $5$ & $4$\\
\bottomrule\end{tabular}\end{center}
\end{exercice}\end{blocLivre}
\begin{blocLivre}{ex-10-7}\begin{exercice}[2]{Une médiane paire}
Pour $1;3;6;8;10;20$, donner la médiane retenue dans ce livre et la moyenne.
\end{exercice}\end{blocLivre}
\begin{blocLivre}{ex-10-8}\begin{exercice}[2]{Un disque statistique}
Sur 40 réponses : 10 choisissent A, 18 choisissent B et 12 choisissent C. Calculer les angles du diagramme circulaire.
\label{fig-ex-10-8}\figureGeo{stat-disque}
\end{exercice}\end{blocLivre}
\begin{blocLivre}{ex-10-9}\begin{exercice}[2]{Un histogramme}
Les durées sont regroupées en classes de 10 minutes : $[0;10[$, $[10;20[$, $[20;30[$, $[30;40[$. Lire les effectifs sur la figure et la classe contenant les deux rangs centraux.
\label{fig-ex-10-9}\figureGeo{stat-histogramme}
\end{exercice}\end{blocLivre}
\begin{blocLivre}{ex-10-10}\begin{exercice}[2]{Une formule de tableur}
Les valeurs 4, 7, 12 sont en A2:A4 et leurs effectifs 3, 5, 2 en B2:B4. Écrire une formule donnant la somme pondérée puis la moyenne.
\end{exercice}\end{blocLivre}

\clearpage\rubrique{Problèmes}
\begin{blocLivre}{ex-10-11}\begin{exercice}[3]{Réunir deux groupes}
Un groupe de 12 élèves a une moyenne de 13 ; un autre de 18 élèves a une moyenne de 11. Quelle est la moyenne des 30 élèves ?
\end{exercice}\end{blocLivre}
\begin{blocLivre}{ex-10-12}\begin{exercice}[3]{Une note manquante}
Quatre notes sont 8, 12, 14 et $x$. Leur moyenne est 12,5. Déterminer $x$.
\end{exercice}\end{blocLivre}
\begin{blocLivre}{ex-10-13}\begin{exercice}[3]{Même moyenne, mêmes résultats ?}
Comparer les séries A : $9;10;10;11$ et B : $2;10;10;18$. Calculer moyenne, médiane et étendue, puis commenter.
\end{exercice}\end{blocLivre}
\begin{blocLivre}{ex-10-14}\begin{exercice}[3]{Ajouter une observation}
Une série de 9 valeurs a une moyenne de 14. On ajoute 24. Calculer la nouvelle moyenne. Peut-on en déduire la nouvelle médiane ?
\end{exercice}\end{blocLivre}
\begin{blocLivre}{ex-10-15}\begin{exercice}[3]{Un graphique trompeur}
Deux barres représentent 80 et 84 inscriptions. Un axe vertical commence à 78. Comparer la hausse réelle et les hauteurs visibles à partir de 78.
\label{fig-ex-10-15}\figureGeo{stat-axe-tronque}
\end{exercice}\end{blocLivre}

\clearpage
